% ── Německý model: Mitbestimmung, Tarifautonomie a dědictví Hartz reforem ────

\paragraph{Institucionální architektura}\label{par:de_architektura}

Německý model stojí na~třech vzájemně propojených pilířích: ústavně
zakotvené \emph{Tarifautonomii} (autonomie sociálních partnerů od~státního
zásahu), odvětvovém vyjednávání mezi konfederacemi (\textit{DGB} na~straně
zaměstnanců, \textit{BDA} na~straně zaměstnavatelů) a~podnikovém
spolurozhodování (\textit{Mitbestimmung}) skrze závodní rady
(\textit{Betriebsrat}) a~účast zástupců zaměstnanců v~dozorčích
radách.
Stát formálně do~vyjednávání nevstupuje --- výjimkou je~mechanismus
\textit{Allgemeinverbindlicherklärung} (rozšíření \acs{KS} \emph{erga omnes}
ministerským dekretem), který je~však podmíněn souhlasem obou stran v~\textit{Tarifausschuss}
a~v~praxi se~uplatňuje pouze v~několika málo odvětvích. Reformami
\emph{Hartz~I--IV} (2003--2005) byl tradiční model doplněn o~prvky
deregulace agenturní práce, mini-jobs a~aktivační podpory --- s~kontroverzními
sociálními dopady.

\paragraph{Pokrytí \acgen{KV} a~hustota odborů}\label{par:de_pokryti}

Pokrytí \acs{KS} v~Německu dlouhodobě klesá: ze~zhruba \SI{85}{\percent}
v~polovině 90.~let na~současných přibližně \SI{52}{\percent} (data \acs{ICTWSS}
pro~rok~2024, obr.~\ref{fig:eu_pokryti_kv_mapa}); pokles je~nejvýraznější
v~nových spolkových zemích (východní Německo cca \SI{45}{\percent} oproti
\SI{55}{\percent} na~západě).
Hustota odborů paralelně klesá pod \SI{16}{\percent} (obr.~\ref{fig:eu_hustota_mapa}).
Rozpad pokrytí je~strukturálně způsoben odchodem firem ze~zaměstnavatelských
svazů (jev označovaný jako~\textit{Tarifflucht}); \textit{WSI Tarifarchiv}
identifikuje jako~hlavní ohniska poklesu maloobchod, ICT, gastronomii
a~osobní služby. V~reakci na~tento
trend přijala spolková vláda v~roce~2025 zákon o~posílení tarifní vazby
(\textit{Tariftreuegesetz}), který podmíňuje veřejné zakázky dodržováním
\acs{KS}.

\paragraph{Klíčové rysy modelu}\label{par:de_rysy}

Zákonná minimální mzda byla zavedena teprve v~roce~2015
(\textit{Mindestlohngesetz}); do~té doby byl mzdový floor stanovován výlučně
\acs{KS}. \emph{Mitbestimmung} představuje nejhlubší forma podnikové
participace zaměstnanců v~\acs{EU}: ve~firmách nad 2~000 zaměstnanců tvoří
zástupci zaměstnanců polovinu dozorčí rady (\textit{paritätische Mitbestimmung}),
v~menších firmách (500--2~000) jednu třetinu. Empirické studie dokládají, že~firmy
s~paritní účastí vykazují vyšší investiční horizont, nižší fluktuaci
a~stabilnější mzdový vývoj v~období hospodářských šoků.

Hartz reformy (2003--2005) přinesly současně pokles nezaměstnanosti
(z~přibližně \SI{10}{\percent} v~roce~2005 na~\SI{5}{\percent} v~letech~2018--2019)
i~výrazný nárůst nestandardních forem zaměstnání (mini-jobs do~\SI{520}{\eur}
měsíčně bez sociálních odvodů, agenturní práce, krátkodobé úvazky). Hartz~IV
sloučil dávky v~nezaměstnanosti se~sociální podporou do~jediného základního
zabezpečení (\textit{Arbeitslosengeld II}, dnes \textit{Bürgergeld}), což snížilo
náhradovou míru pro~dlouhodobě nezaměstnané a~zvýšilo chudobu pracujících
(\emph{working poor}). Současná
debata Bundestagu se~zabývá kompenzací těchto dopadů skrze posílení \acs{KV}
a~regulaci platformové práce.

Z~pohledu výhledových faktorů ekonomický výzkumný ústav \textit{IAB}
poukazuje na~průmysl~4.0 jako~na~dvojí příležitost: na~jedné straně
automatizace ohrožuje rutinní pracovní místa nízké kvalifikace, na~druhé
straně koordinované \acs{KV} a~\textit{Mitbestimmung} představují institucionální
mechanismus pro~strukturovanou rekvalifikaci a~spravedlivé rozdělení
produktivních zisků.

\paragraph{Implikace pro~\acacc{ČR}}\label{par:de_implikace}

\Acgen{ČR} sdílí s~Německem strukturu ekonomiky orientované na~průmyslový
export, ale postrádá ústavní zakotvení \emph{Tarifautonomie} a~účinný mechanismus
rozšiřování \acs{KSVS}. Nejcennější přenositelnou součástí německého modelu je~institut
podnikových rad (závodních rad) s~právem na~informace, konzultaci a~spolurozhodování
v~personálních otázkách. Rada zaměstnanců dle~\S~281~zákoníku práce~\cite{zakon_zp_2006}
v~\acgen{ČR} formálně existuje, její pravomoci jsou však výrazně užší než
\textit{Betriebsrat} v~Německu. Posílení role rad zaměstnanců a~jejich
zákonného nároku na~konzultaci v~otázkách restrukturalizace by~zlepšilo
adaptační kapacitu trhu práce vůči vlnám automatizace, aniž by~vyžadovalo
změnu \acs{KV} systému jako~celku. Současně varuje zkušenost s~Hartz reformami,
že~deregulace okrajových forem zaměstnání bez paralelního posílení \acs{KV}
vede k~erozi mzdové struktury a~k~růstu chudoby pracujících.
